\section{Experiments}
\label{sec:experiments}
We evaluate our method along four axes:
\textit{(i)} dynamic scene reconstruction quality on Waymo Open Dataset,
covering temporal interpolation and extrapolation
(\S~\ref{sec:temporal_evaluation});
\textit{(ii)} off-trajectory view synthesis under lateral
camera shifts (\S~\ref{sec:spatial_evaluation});
\textit{(iii)} cross-dataset transfer from Waymo to nuScenes
(\S~\ref{sec:generalization}); and
\textit{(iv)} ablations of the main design choices
(\S~\ref{sec:ablations}).

\subsection{Experimental Setup}
\label{sec:experimental_setup}

\paragraph{Datasets and protocol.}
We use the Waymo Open Dataset with 798 training and 202 validation sequences. Following STORM, each validation sequence is divided into non-overlapping 2-second clips of 20 frames at 10 \, Hz~\cite{yang2025storm}. The FRONT, FRONT\_LEFT, and FRONT\_RIGHT cameras provide observations at zero-based frame indices $\{0,5,10,15\}$. The remaining 16 frames are used for evaluation: 12 frames for temporal interpolation and four frames for prediction beyond the last observation. Predictions and reference images are resized to $240\times160$ pixels for evaluation~\cite{huang2026instantnurec}.

\paragraph{Baselines.}
We compare with STORM, DGGT, StreetForward~\cite{yang2025storm,chen2025dggt,yu2026streetforward}. Methods use the same context timestamps and target views, with model-specific preprocessing. Camera parameters, additional geometric inputs, and rendering adapters are documented with the implementation settings. For spatial extrapolation, we compare Gaussian renderings before diffusion refinement, following PointForward~\cite{chi2026pointforward}.

\paragraph{Metrics.}
We report PSNR and SSIM over the full image and shared annotated regions. Depth quality is measured using AbsRel and $\delta_1$ at valid projected LiDAR samples after per-frame scale alignment~\cite{huang2026instantnurec}. Scores are averaged over frames within each clip and then over clips within each sequence; sequences receive equal weight. Novel-trajectory renderings are evaluated using FID and FVD against a fixed real-image and video reference population.

\paragraph{Implementation details.}
DriveCompose reconstructs category-specific actor assets from motion-aligned observations and composes them with a static background. Reconstruction time includes preprocessing, asset generation, and scene composition. Rendering time is measured separately, and asset reuse is evaluated across repeated scene queries.
% Fill in verified DriveCompose training schedule, resolution, hardware and checkpoint.
% Do not copy baseline training configurations into DriveCompose implementation details.

\subsection{Dynamic Reconstruction and Temporal Prediction}
\label{sec:temporal_evaluation}

We evaluate reconstruction at unobserved times along the recorded camera trajectory. Table~\ref{tab:regions_current} reports image and depth quality. The regional comparison focuses on actor appearance, and synchronized visualizations examine object position, shape, and occlusion changes.
% Add measured DriveCompose findings here after replacing the development numbers.
\begin{table}[!htbp]
\centering
\footnotesize
\setlength{\tabcolsep}{1.3pt}
\renewcommand{\arraystretch}{1.08}
\setlength{\belowcaptionskip}{5pt}
\caption{Dynamic scene reconstruction on Waymo. Image quality is reported over the full image and annotated regions; depth is evaluated at valid LiDAR samples.}
\label{tab:regions_current}
\begin{tabular}{@{}lrrrrrr@{}}
\toprule
& \multicolumn{2}{c}{Full image} & \multicolumn{2}{c}{\shortstack{Dynamic Region}} & \multicolumn{2}{c}{Depth} \\
\cmidrule(lr){2-3}\cmidrule(lr){4-5}\cmidrule(lr){6-7}
Method & PSNR $\uparrow$ & SSIM $\uparrow$ & PSNR $\uparrow$ & SSIM $\uparrow$ & AbsRel $\downarrow$ & $\delta_1$ $\uparrow$ \\
\midrule
STORM & 26.191 & 0.833 & 22.083 & 0.654 & 0.129 & 0.867 \\
DGGT & 26.056 & 0.809 & 20.809 & 0.603 & 0.168 & 0.779 \\
StreetForward & 25.256 & 0.814 & 21.844 & 0.636 & 0.122 & 0.860 \\
\midrule
DriveCompose & -- & -- & -- & -- & -- & -- \\
\bottomrule
\end{tabular}
\end{table}

\paragraph{Temporal interpolation and prediction.}
Table~\ref{tab:temporal_current} separates temporal interpolation from prediction at 0.1--0.4 seconds beyond the last context frame.

\begin{table}[!htbp]
\centering
\footnotesize
\setlength{\tabcolsep}{1.3pt}
\renewcommand{\arraystretch}{1.08}
\setlength{\belowcaptionskip}{5pt}
\caption{Temporal interpolation and extrapolation on Waymo. PSNR and SSIM are reported separately for interpolation and extrapolation; context frames are excluded.}
\label{tab:temporal_current}
\begin{tabular}{@{}lrrrrrr@{}}
\toprule
& \multicolumn{2}{c}{Interp.} & \multicolumn{2}{c}{Extrap.} & \multicolumn{2}{c}{All targets} \\
\cmidrule(lr){2-3}\cmidrule(lr){4-5}\cmidrule(lr){6-7}
Method & PSNR $\uparrow$ & SSIM $\uparrow$ & PSNR $\uparrow$ & SSIM $\uparrow$ & PSNR $\uparrow$ & SSIM $\uparrow$ \\
\midrule
STORM & 26.427 & 0.842 & 25.481 & 0.807 & 26.191 & 0.833 \\
DGGT & 26.265 & 0.817 & 25.429 & 0.788 & 26.056 & 0.809 \\
StreetForward & 25.653 & 0.827 & 24.065 & 0.776 & 25.256 & 0.814 \\
\midrule
DriveCompose & -- & -- & -- & -- & -- & -- \\
\bottomrule
\end{tabular}
\end{table}

\subsection{Off-Trajectory View Synthesis}
\label{sec:spatial_evaluation}

We shift the camera trajectory laterally by 1, 2, and 4 meters to the left and right, keeping camera orientation, intrinsics, and timestamps fixed. This evaluation tests the stability of scene structure and the rendering of newly visible surfaces as the viewpoint moves away from the observations~\cite{zhu2025worldsplat,chi2026pointforward}.

Table~\ref{tab:spatial_current} reports FID and FVD
at each displacement magnitude, averaged over separately
evaluated left and right shifts. Both metrics use the
same real reference population across methods.
Qualitative comparisons use matched scenes, timestamps,
and camera offsets to examine vehicle shape, road layout,
and occlusion boundaries.

\begin{table}[t]
\centering
\footnotesize
\setlength{\tabcolsep}{2.5pt}
\renewcommand{\arraystretch}{1.08}
\caption{Novel-view synthesis under lateral camera shifts.
FID and FVD are averaged over left and right shifts
at each displacement magnitude. Lower is better.}
\label{tab:spatial_current}
\resizebox{\columnwidth}{!}{%
\begin{tabular}{@{}l|rr|rr|rr@{}}
\toprule
Method
& \multicolumn{2}{c|}{Shift $\pm 1$ m}
& \multicolumn{2}{c|}{Shift $\pm 2$ m}
& \multicolumn{2}{c}{Shift $\pm 4$ m} \\
\cmidrule(lr){2-3}
\cmidrule(lr){4-5}
\cmidrule(lr){6-7}
& FID $\downarrow$ & FVD $\downarrow$
& FID $\downarrow$ & FVD $\downarrow$
& FID $\downarrow$ & FVD $\downarrow$ \\
\midrule
STORM
& 61.97 & 722.27
& 83.72 & 1029.36
& 132.57 & 1704.16 \\
DGGT
& 73.99 & 984.00
& 93.36 & 1197.03
& 128.01 & 1661.70 \\
StreetForward
& 84.30 & 878.95
& 103.31 & 1095.02
& 123.58 & 1578.63 \\
\midrule
\rowcolor{gray!10}
DriveCompose
& -- & --
& -- & --
& -- & -- \\
\bottomrule
\end{tabular}%
}
\end{table}

% Avoid stretching paragraph gaps beside the taller single-column table.
\raggedbottom
\subsection{Cross-Dataset Generalization}
\label{sec:generalization}

We evaluate transfer from Waymo to nuScenes under zero-shot and fine-tuning settings. Zero-shot evaluation directly applies the Waymo-trained checkpoint without target-domain updates. Fine-tuning adapts the same initialization using the nuScenes training split and evaluates it on the held-out validation split. Both settings use the same target evaluation scenes, context sampling, and metrics.

Table~\ref{tab:generalization} reports the two settings separately. The zero-shot comparison assesses transfer to new environments, while the fine-tuning comparison measures adaptation to the target domain. Full-image and annotated-region scores distinguish overall scene quality from actor reconstruction.
% Planned protocol: all cells below are unmeasured. Fill the target split,
% fine-tuning data budget, update count and camera mapping before evaluation.
% Span both columns so the generalization comparison is not stranded in the
% right column beside an empty float area.

\begin{table}[t]
\centering
\footnotesize
\setlength{\tabcolsep}{3.5pt}
\renewcommand{\arraystretch}{1.08}
\caption{Cross-dataset generalization and target-domain adaptation.
Zero-shot models are evaluated directly on the target datasets, while
trained models are fine-tuned on the target domain. Higher PSNR and SSIM,
and lower LPIPS are better.}
\label{tab:generalization}
\begin{tabular*}{\columnwidth}{
    @{\extracolsep{\fill}}l
    ccc@{\hskip 5pt}|@{\hskip 5pt}ccc
}
\toprule
& \multicolumn{3}{c}{\textbf{Full image}}
& \multicolumn{3}{c}{\textbf{Dynamic regions}} \\
\cmidrule(lr){2-4}\cmidrule(lr){5-7}
\textbf{Method}
& PSNR $\uparrow$ & SSIM $\uparrow$ & LPIPS $\downarrow$
& PSNR $\uparrow$ & SSIM $\uparrow$ & LPIPS $\downarrow$ \\
\midrule

\multicolumn{7}{l}{\textbf{Zero-shot}} \\
STORM
& -- & -- & --
& -- & -- & -- \\
DGGT
& -- & -- & --
& -- & -- & -- \\
StreetForward
& -- & -- & --
& -- & -- & -- \\
\rowcolor{gray!12}
\textbf{DriveCompose}
& \textbf{--} & \textbf{--} & \textbf{--}
& \textbf{--} & \textbf{--} & \textbf{--} \\

\midrule
\multicolumn{7}{l}{\textbf{Trained}} \\
STORM
& -- & -- & --
& -- & -- & -- \\
\rowcolor{gray!12}
\textbf{DriveCompose}
& \textbf{--} & \textbf{--} & \textbf{--}
& \textbf{--} & \textbf{--} & \textbf{--} \\

\bottomrule
\end{tabular*}
\end{table}

\subsection{Ablation protocol}
\label{sec:ablation-protocol}
\label{sec:ablations}

\paragraph{Inputs and controls.}
Paired evaluation of hidden surfaces is defined on independent complete-reference scenes; Waymo logs provide recorded-view evaluation.
Each 29-timestamp window uses three synchronized cameras: front, front-left, and front-right.
The nested input sets are $\mathcal T_2=\{0,28\}$, $\mathcal T_4=\{0,8,20,28\}$, and $\mathcal T_8=\{0,4,\ldots,28\}$, with shared target timestamps $\{2,6,\ldots,26\}$.
Target images are excluded from evidence and query construction; target camera and actor states are supplied equally to every variant.
The recovery module and codec are frozen, with matched query layouts, data exposure, and optimization budgets at each input setting.
Training samples $T\in\{2,4,8\}$ uniformly; completion uses 24 Euler steps.
The four learned variants use training seeds $\{0,1,2\}$ and world seeds $\{7,17,27\}$; A1 is deterministic.
Hyperparameters and acceptance thresholds are selected on an internal development split before evaluation on scene- and identity-disjoint data.

\paragraph{Metrics.}
New-MAE measures linear-RGB error on independently identified surfaces unseen in every input; NVS-MAE covers all held-out target pixels.
Across input settings, New-MAE uses the common region unseen even under $\mathcal T_8$.
Obs-Reg averages the positive input-view regression in Eq.~\ref{eq:regression}.
Cov.\ is the non-sky target fraction with support coverage $>0.5$, opacity $>0.1$, and finite positive depth.
Uncovered pixels receive the same zero-RGB fill for error measurement.
Pass is the fraction of all requested worlds accepted by the fixed checker, including construction failures in its denominator.
Image metrics include every renderable candidate, including rejected worlds; non-renderable counts are recorded separately.
Metrics are averaged within each scene and then across scenes and training seeds.

Figure~\ref{fig:ablation-reuse} reserves matched crops for completion and observation preservation, followed by three state interventions from the same world.

\begin{figure*}[t]
  \centering
  \begingroup
  \newcommand{\dcablationcell}[3]{%
    \begin{minipage}[t]{#1}
      \centering
      \parbox[c][2.2\baselineskip][c]{\linewidth}{\centering\small #3}\par
      \fbox{\rule{0pt}{#2}\rule{.90\linewidth}{0pt}}
    \end{minipage}}
  {\small\textbf{(a) Unseen-surface completion}}\par\smallskip
  \dcablationcell{.158\textwidth}{12mm}{GT}\hfill
  \dcablationcell{.158\textwidth}{12mm}{No completion\\(A1)}\hfill
  \dcablationcell{.158\textwidth}{12mm}{Independent\\(A2)}\hfill
  \dcablationcell{.158\textwidth}{12mm}{No world coord.\\(A3)}\hfill
  \dcablationcell{.158\textwidth}{12mm}{Soft observation\\(A7)}\hfill
  \dcablationcell{.158\textwidth}{12mm}{Full\\(A0)}
  \par\medskip
  {\small\textbf{(b) Input-view preservation}}\par\smallskip
  \dcablationcell{.158\textwidth}{12mm}{GT / Input}\hfill
  \dcablationcell{.158\textwidth}{12mm}{No completion\\(A1)}\hfill
  \dcablationcell{.158\textwidth}{12mm}{Independent\\(A2)}\hfill
  \dcablationcell{.158\textwidth}{12mm}{No world coord.\\(A3)}\hfill
  \dcablationcell{.158\textwidth}{12mm}{Soft observation\\(A7)}\hfill
  \dcablationcell{.158\textwidth}{12mm}{Full\\(A0)}
  \par\medskip
  {\small\textbf{(c) One fixed world, different states (A0)}}\par\smallskip
  \dcablationcell{.235\textwidth}{20mm}{Original $s_0$}\hfill
  \dcablationcell{.235\textwidth}{20mm}{Camera $+1\,\mathrm m$}\hfill
  \dcablationcell{.235\textwidth}{20mm}{Vehicle $+30^\circ$}\hfill
  \dcablationcell{.235\textwidth}{20mm}{Human pose $H_1$}
  \endgroup
  \caption{\textbf{Ablation and reuse layout.} (a) The same newly exposed region across five variants and an independent reference. (b) Matched crops in an original input view. (c) Independent camera, vehicle, and human-state changes from one A0 world with seed 7; the other states and lighting are fixed. All image slots are intentionally empty.}
  \label{fig:ablation-reuse}
\end{figure*}

\subsection{Ablation studies}
\label{sec:ablation-studies}

\paragraph{Completion and scene context.}
Table~\ref{tab:ablation-components} compares Full (A0) with three controls.
No completion (A1) disables the opacity of every non-observed, non-sky query.
Independent (A2) isolates component memories at every layer while sharing parameters and retaining the same global log context.
No world coord.\ (A3) zeros only the explicit reference-world coordinate channels in the completion network; canonical geometry and rendering transforms remain unchanged.
The learned controls are retrained through Stages 3--4; memory-related compute differences are recorded.

\paragraph{Observation preservation.}
Soft observation (A7) allows observed codes to evolve instead of hard resetting them.
Its observed flow path is $(1-\tau)\mathbf z^{\rm obs}+\tau\mathbf z_1$, with target velocity $\mathbf z_1-\mathbf z^{\rm obs}$.
Observed and unknown flow losses are separately normalized and equally weighted.
Rendering adaptation adds a mean-squared penalty on observed-code changes, with its weight selected from $\{0.1,1,10\}$ on development data.
Observed conditioning, free-space constraints, the codec, and the checker remain fixed.
This comparison examines the tradeoff between completion and input-view preservation.

\begin{table}[t]
  \centering
  \caption{\textbf{Component ablations.} All variants use $T=8$, a shared recovery module and codec, and the same evaluation masks. Quality is measured before acceptance filtering; Pass includes all requested worlds. A1 requires no completion training. Numerical cells are left blank.}
  \label{tab:ablation-components}
  \small
  \setlength{\tabcolsep}{3pt}
  \begin{tabular}{@{}lrrrr@{}}
    \toprule
    Variant & \shortstack{New-MAE\\$\downarrow$}
      & \shortstack{Obs-Reg\\$\downarrow$}
      & \shortstack{Cov.\\(\%) $\uparrow$}
      & \shortstack{Pass\\(\%) $\uparrow$}\\
    \midrule
    Full (A0)            & & & & \\
    No completion (A1)   & & & & \\
    Independent (A2)     & & & & \\
    No world coord. (A3) & & & & \\
    Soft obs. (A7)       & & & & \\
    \bottomrule
  \end{tabular}
\end{table}

\paragraph{Input sparsity.}
Table~\ref{tab:ablation-inputs} reuses the A0/A2 checkpoints at $T=2,4,8$, with 6/12/24 input images.
The time span, target views, and common unseen-region mask remain fixed; evidence and queries are rebuilt from each input subset.
No model is retrained for this comparison.

\begin{table}[t]
  \centering
  \caption{\textbf{Input-timestamp sensitivity.} Three cameras and a fixed time span are shared across the nested input sets. The same checkpoints and target images are used at every $T$; New-MAE uses the common unseen region. Numerical cells are left blank.}
  \label{tab:ablation-inputs}
  \small
  \setlength{\tabcolsep}{3pt}
  \begin{tabular}{@{}lr rrr@{}}
    \toprule
    Variant & $T$ & \shortstack{NVS-MAE\\$\downarrow$}
      & \shortstack{New-MAE\\$\downarrow$}
      & \shortstack{Cov.\\(\%) $\uparrow$}\\
    \midrule
    Full (A0)        & 2 & & & \\
    Independent (A2) & 2 & & & \\
    Full (A0)        & 4 & & & \\
    Independent (A2) & 4 & & & \\
    Full (A0)        & 8 & & & \\
    Independent (A2) & 8 & & & \\
    \bottomrule
  \end{tabular}
\end{table}

\subsection{State-driven reuse}
\label{sec:ablation-reuse}


At timestamp 14 and world seed 7, the interventions independently translate the camera by $1\,\mathrm m$ along its local $+x$ axis, rotate a vehicle by $30^\circ$ about its root and the world vertical axis, or apply a held-out valid human pose $H_1$ with fixed root placement.
Lighting and canonical assets remain fixed, without resampling completion.

% Flush experiment floats before Discussion and Conclusion.
\FloatBarrier
